% Gram matrix rigidity: what the Lean formalization proves, and how it relates to
% the literature.  Expert-readable companion to
% ForMathlib/Analysis/InnerProductSpace/GramMatrix.lean, so the statement can be
% checked against the source references without reading Lean.
%
% Independently written; draws on the distilled digest
%   ForMathlib/prose/Horn-Johnson-2013-Gram-core-arguments.tex

\documentclass[11pt]{article}

\usepackage[T1]{fontenc}
\usepackage{lmodern}
\usepackage{microtype}
\usepackage[margin=1.1in]{geometry}
\usepackage{amsmath,amssymb,amsthm,mathtools}
\usepackage{booktabs}
\usepackage[colorlinks=true, urlcolor=blue, citecolor=blue, linkcolor=black]{hyperref}

\newcommand{\M}{M}
\newcommand{\norm}[1]{\left\lVert #1\right\rVert}
\newcommand{\inner}[2]{\langle #1,\,#2\rangle}
\newcommand{\adj}{^{*}}
\newcommand{\KK}{\mathbb K}
\newcommand{\gram}{G}
\DeclareMathOperator{\rank}{rank}
\DeclareMathOperator{\spann}{span}
\DeclareMathOperator{\ran}{ran}

\newtheorem{theorem}{Theorem}
\newtheorem{corollary}[theorem]{Corollary}
\newtheorem{lemma}[theorem]{Lemma}
\theoremstyle{remark}
\newtheorem*{remark}{Remark}

\title{Gram matrix rigidity: what the Lean formalization proves,\\
and how it relates to Horn--Johnson (2013) and finite frames}

% NOTE: if you are a different model making a substantial addition and are not
% yet credited, add yourself to the author list above.
\author{Claude Opus 4.8 \and Jonathan Crall}

\date{\today}

\begin{document}
\maketitle

\begin{abstract}
This note states, in ordinary mathematical language, the Gram-matrix rigidity
theorem formalized in Lean, and locates it against the source references.  Two
families of vectors with the same Gram matrix, that is with equal pairwise inner
products, are related by a linear isometry: on their spans in full generality,
and on a finite-dimensional ambient space by a linear isometry equivalence.  This
is the coordinate-free form of the classical fact that a factorization
$A\adj A=B\adj B$ determines $A$ up to a left isometry (Horn and Johnson,
Theorem 7.3.11), and of the frame-theoretic statement that finite frames are
unitarily equivalent if and only if their Gram matrices coincide.  The formalized
proof replaces the singular value decomposition (SVD) with an isometric first
isomorphism theorem, and needs no finite-dimensionality for the span-level
statement.  A closing section records what is left unformalized.
\end{abstract}

\vspace{1.5em}
% NOTE: Update permalink if underlying proof changes, so this always is up to date wrt to the referenced lean file.
\noindent Companion note to \href{https://github.com/AIQ-Kitware/aiq-dkps-formalization/blob/37d857e52be20d145cb90f29e8fe5d2a0928c3d1/ForMathlib/Analysis/InnerProductSpace/GramMatrix.lean}{ForMathlib/Analysis/InnerProductSpace/GramMatrix.lean}, part of a larger formalization effort at \href{https://github.com/AIQ-Kitware/aiq-dkps-formalization}{github.com/AIQ-Kitware/aiq-dkps-formalization}.

\section{Setup and notation}

Let $E$ and $F$ be inner product spaces over $\KK\in\{\mathbb R,\mathbb C\}$, and
let $\varphi=(\varphi_i)_{i\in I}$ in $E$ and $\psi=(\psi_i)_{i\in I}$ in $F$ be
two families of vectors sharing an index set $I$.  Their Gram matrices are
$\gram(\varphi)=[\inner{\varphi_i}{\varphi_j}]$ and
$\gram(\psi)=[\inner{\psi_i}{\psi_j}]$.  The families have \emph{equal Gram data}
when
\begin{equation}\label{eq:equalgram}
  \inner{\varphi_i}{\varphi_j}=\inner{\psi_i}{\psi_j}\qquad\text{for all }i,j,
\end{equation}
equivalently $\gram(\varphi)=\gram(\psi)$.  A \emph{linear isometry} is a linear
map preserving the inner product, and a \emph{linear isometry equivalence} is a
bijective one; over $\KK$ these are the (real) orthogonal and (complex) unitary
maps.  Write $\spann\varphi=\spann\{\varphi_i:i\in I\}$.

In matrix terms a family in $\KK^{p}$ is the columns of a matrix $A\in\M_{p,n}$,
its Gram matrix is $A\adj A$, and \eqref{eq:equalgram} is $A\adj A=B\adj B$.  A
positive semidefinite (PSD) matrix is a Hermitian $M$ with $x\adj Mx\ge 0$; it is
positive definite (PD) when the inequality is strict for $x\ne 0$.

\section{The formalized theorems}

Two blocks are formalized: the characterization of Gram matrices, and the
rigidity.

\begin{theorem}[Gram characterization]\label{thm:char}
$\gram(\varphi)$ is Hermitian and positive semidefinite, and it is positive
definite if and only if the family $\varphi$ is linearly independent.
\end{theorem}

\begin{theorem}[Gram rigidity]\label{thm:rigid}
Suppose $\varphi$ and $\psi$ have equal Gram data \eqref{eq:equalgram}.
\begin{enumerate}
\item[(a)] There is a linear isometry equivalence
  $\spann\varphi\xrightarrow{\ \sim\ }\spann\psi$ sending each $\varphi_i$ to
  $\psi_i$.  No finiteness is assumed, and $E,F$ may differ.
\item[(b)] If $F=E$ is finite-dimensional, this extends to a linear isometry
  equivalence $W$ of $E$ with $W\varphi_i=\psi_i$ for all $i$.
\end{enumerate}
Consequently, for families in a finite-dimensional $E$,
\[
  \gram(\varphi)=\gram(\psi)
  \quad\Longleftrightarrow\quad
  \exists\,W\ \text{a linear isometry equivalence of }E\ \text{with }
  W\varphi_i=\psi_i .
\]
\end{theorem}

Part (a) is the fundamental statement; (b) and the equivalence follow from it by
extending the span isometry to the ambient space.  The aligning map in (a) is
determined on $\spann\varphi$, so it is unique.

\section{The proof mechanism}

The proof is coordinate-free and uses no SVD.  Encode each family by its
linear-combination map $L_\varphi:\KK^{(I)}\to E$, $c\mapsto\sum_i c_i\varphi_i$,
whose range is $\spann\varphi$.  Equal Gram data \eqref{eq:equalgram} makes the
two maps carry the same pullback inner product,
$\inner{L_\varphi c}{L_\varphi c'}=\inner{L_\psi c}{L_\psi c'}$, by expanding
bilinearly over the entries; in particular $L_\varphi c=0$ if and only if
$L_\psi c=0$, since a vector vanishes exactly when its self-inner-product does, so
$\ker L_\varphi=\ker L_\psi$.  An isometric refinement of the first isomorphism
theorem then applies: the assignment $L_\varphi c\mapsto L_\psi c$ is a
well-defined linear isometry equivalence $\ran L_\varphi\to\ran L_\psi$, that is
$\spann\varphi\to\spann\psi$, sending $\varphi_i$ to $\psi_i$.  In finite
dimension it extends to the ambient space.  The engine is stated once, as a fact
about a pair of linear maps $S,T$ with $\inner{Sx}{Sy}=\inner{Tx}{Ty}$: they have
equal kernels and canonically isometric ranges via $Sx\mapsto Tx$.

\section{Relation to Horn--Johnson (2013)}

Horn and Johnson~\cite{HornJohnson2013} record both blocks.  The characterization
is Theorem 7.2.10: a matrix is a Gram matrix if and only if it is positive
semidefinite, it is positive definite if and only if the vectors are linearly
independent, and $\rank\gram=\dim\spann\{v_i\}$; the proof is the identity
$x\adj\gram x=\bigl\|\sum_i x_i v_i\bigr\|^2$, which is the same computation
behind Theorem~\ref{thm:char}.  The rigidity is Theorem 7.3.11: for $p\le q$,
$A\in\M_{p,n}$, and $B\in\M_{q,n}$, one has $A\adj A=B\adj B$ if and only if
$B=VA$ for some $V\in\M_{q,p}$ with orthonormal columns (an isometry).  The square
case is Problem 2.6.P10, and a positive-definite special case, $C\adj C=A$ giving
$C=VA^{1/2}$, is a \S7.2 exercise.

Their proof of 7.3.11 uses the SVD.  Writing $A=V_1\Sigma_r W_1\adj$ and
$B=V_2\Sigma_r W_1\adj$, the two matrices share the right factor $\Sigma_r W_1\adj$
because $A\adj A=W_1\Sigma_r^2 W_1\adj=B\adj B$, so they differ only in the
left-singular frames $V_1,V_2$, and $V_2=VV_1$ for an isometry $V$ obtained by
extending $V_1$ to a unitary.  The formalized proof is this argument with the
shared right factor replaced by the common coimage $\KK^{(I)}/\ker L_\varphi$: the
quotient plays the role of $\Sigma_r W_1\adj$, and extending the span isometry to
the ambient space plays the role of extending $V_1$ to a unitary.  Two differences
of scope: the formalized part~(a) needs no finiteness and allows $E\ne F$, whereas
7.3.11 is finite matrices with $p\le q$; and the construction is by the quotient
rather than the SVD.

\section{Relation to finite frames}

In frame theory a family $(\varphi_i)$ in $\KK^{d}$ is a finite frame, and
Theorem~\ref{thm:rigid}(b) reads: two finite frames are unitarily equivalent,
$\varphi_i\mapsto U\varphi_i=\psi_i$ for one unitary $U$, if and only if they have
the same Gram matrix.  Chien and Waldron~\cite{ChienWaldron2016} study the finer
\emph{projective} unitary equivalence, where the frames agree up to unimodular
scalars, $\varphi_i\mapsto z_i\,U\varphi_i$ with $|z_i|=1$.  Projective equivalence
is not detected by the Gram matrix alone; it is characterized by a family of
higher frame invariants, products of inner products taken around cycles of
indices.  That refinement is not formalized here.

\section{What is not formalized}

The formalization covers parts (a),(b) of the characterization 7.2.10 and the
isometric rigidity.  The remaining pieces, and what each needs, are as follows.

\begin{itemize}
\item \emph{The rank identity} $\rank\gram(\varphi)=\dim\spann\{\varphi_i\}$
  (7.2.10(c)).  It rests on the rank-principal structure of a positive
  semidefinite matrix.
\item \emph{Realization in the reverse direction:} that every positive
  semidefinite $G$ is the Gram matrix of an explicit family, the columns of
  $G^{1/2}$, with a rank-$d$ refinement.  This is a separate development, driven
  by the positive semidefinite square root.
\item \emph{The explicit isometry via the SVD.}  The formalized $W$ is obtained
  abstractly, through the coimage and an extension step; the singular value
  decomposition that would present it in coordinates is not built.
\item \emph{Uniqueness as a stand-alone theorem.}  The aligning map is determined
  on $\spann\varphi$, so uniqueness holds, but it is stated as a remark rather
  than formalized separately.
\item \emph{Projective unitary equivalence} of finite frames up to unimodular
  scalars (Chien--Waldron), characterized by higher frame invariants rather than
  by the Gram matrix.
\item \emph{The companion factorizations} (QR, Cholesky; 7.2.9) as other
  realizations of a positive semidefinite matrix.
\end{itemize}

\section{Dictionary to the Lean source}

The characterization lemmas are in Mathlib's
\texttt{Analysis/InnerProductSpace/GramMatrix.lean}; the rigidity declarations are
the \texttt{ForMathlib} additions staged for the same file.

\begin{center}
\footnotesize
\setlength{\tabcolsep}{4pt}
\begin{tabular}{@{}p{0.46\linewidth}l@{}}
\toprule
Statement here & Lean name \\
\midrule
Gram matrix $\gram(\varphi)=[\inner{\varphi_i}{\varphi_j}]$
  & \texttt{Matrix.gram} \\
Thm~\ref{thm:char}, positive semidefinite (7.2.10a)
  & \texttt{Matrix.posSemidef\_gram} \\
Thm~\ref{thm:char}, positive definite $\iff$ independent (7.2.10b)
  & \texttt{Matrix.posDef\_gram\_iff\_linearIndependent} \\
Factorization $\gram=m\adj m$ (7.2.7)
  & \texttt{Matrix.gram\_eq\_conjTranspose\_mul} \\
Engine: isometric first isomorphism theorem
  & \texttt{LinearMap.rangeEquivOfInnerEq} \\
Thm~\ref{thm:rigid}(a), span isometry
  & \texttt{linearIsometryEquivSpanOfInnerEq} \\
Thm~\ref{thm:rigid}(b), ambient isometry equivalence
  & \texttt{exists\_linearIsometryEquiv\_map\_eq\_of\_inner\_eq} \\
Thm~\ref{thm:rigid}, the $\gram(\varphi)=\gram(\psi)$ iff
  & \texttt{Matrix.gram\_eq\_gram\_iff\_\dots\_map\_eq} \\
\bottomrule
\end{tabular}
\end{center}

The span isometry of Theorem~\ref{thm:rigid}(a) is
\texttt{linearIsometryEquivSpanOfInnerEq}, built from the engine
\texttt{LinearMap.rangeEquivOfInnerEq} applied to the two linear-combination maps;
the ambient extension uses \texttt{LinearIsometry.extend}.

\begin{thebibliography}{9}
\bibitem{HornJohnson2013}
Roger A.\ Horn and Charles R.\ Johnson.
\newblock \href{https://doi.org/10.1017/CBO9781139020411}{\emph{Matrix Analysis}}, 2nd ed.
\newblock Cambridge University Press, 2013.
\newblock Sections 7.2--7.3; Problem 2.6.P10.

\bibitem{ChienWaldron2016}
Tuan-Yow Chien and Shayne Waldron.
\newblock \href{https://doi.org/10.1137/15M1042140}{A characterization of projective unitary equivalence of finite frames and applications}.
\newblock \emph{SIAM J.\ Discrete Math.}\ \textbf{30} (2016), no.\ 2, 976--994.
\newblock arXiv:\href{https://arxiv.org/abs/1312.5393}{1312.5393}.
\end{thebibliography}

\end{document}
